\emph{Översikt}
Övningen är sju uppgifter att lösa själv, med ett lösningsförslag till
var och en.  De första handlar om att välja behållare: en mängd svarar på
vilka olika värden som förekommer, en uppslagslista på hur många gånger
vart och ett förekommer.  Sedan går vi över till klasserna: vi utökar
bråkklassen så att bråk går att jämföra och dividera --- och upptäcker på
vägen vad som krävs av ett objekt som ska duga som nyckel --- och bygger
en bank vars konton ligger i en behållare inuti klassen, med en stack som
minns vad banken gjort så att det går att ångra.  De två sista
uppgifterna är fördjupning: en radbrytande ko och två moduler som
tillsammans blir ett program.  Varje uppgift går att lösa på fler sätt än
ett, och lösningsförslagen säger var valen finns.

\emph{Lärandemål}
Efter övningen ska studenten kunna:
% Lo-koderna i kommentarerna nedan är veckosidans egna, och hör inte i
% studentens text.
\begin{restatable}{lo}{OvningContainersLOChoose}%
  \label{OvningContainersLOChoose}%
  Välja behållare --- lista, tupel, mängd, uppslagslista, stack eller
  kö --- utifrån vilka operationer programmet behöver, och motivera
  valet.
\end{restatable}% Lo09, Lo02
\begin{restatable}{lo}{OvningContainersLOSet}%
  \label{OvningContainersLOSet}%
  Använda en mängd för att svara på vilka olika värden som förekommer,
  och förklara varför den varken har ordning eller dubbletter.
\end{restatable}% Lo09
\begin{restatable}{lo}{OvningContainersLODict}%
  \label{OvningContainersLODict}%
  Bygga upp en uppslagslista medan programmet körs, hantera att en
  nyckel saknas, och gå igenom både nycklar och värden.
\end{restatable}% Lo09, Lo04
\begin{restatable}{lo}{OvningContainersLOOperators}%
  \label{OvningContainersLOOperators}%
  Implementera operatoröverlagring med dundermetoder för jämförelse och
  division, och avgöra när en omvänd operator kan återanvända den
  vanliga.
\end{restatable}% Lo09
\begin{restatable}{lo}{OvningContainersLOKey}%
  \label{OvningContainersLOKey}%
  Förklara vad som krävs av ett objekt av en egen klass för att det ska
  kunna vara nyckel i en uppslagslista eller element i en mängd.
\end{restatable}% Lo09
\begin{restatable}{lo}{OvningContainersLOContainer}%
  \label{OvningContainersLOContainer}%
  Konstruera program där en klass har en behållare av andra objekt som
  attribut, och där metoderna söker i och uppdaterar behållaren.
\end{restatable}% Lo09, Lo02
\begin{restatable}{lo}{OvningContainersLODocs}%
  \label{OvningContainersLODocs}%
  Hitta och läsa dokumentationen för Pythons behållare och därifrån
  använda en behållare eller metod som inte gåtts igenom på
  föreläsningen.
\end{restatable}% Lo09

\emph{Förkunskaper}
Studenten bör ha tagit del av föreläsningarna \emph{Behållare:
Uppslagslistor}, som ger uppslagslistan och hur en saknad nyckel
hanteras, \emph{Behållare: Mängder, stackar och köer}, som ger mängden,
stacken och kön samt kriteriet för att välja behållare, \emph{Behållare:
Ett gissningsspel}, som visar hur delarna sätts ihop till ett program och
hur en körning görs upprepbar, \emph{Operatoröverlagring}, som bygger den
bråkklass två av uppgifterna utgår från, och \emph{Praktiska
tillämpningar av klasser}, som visar en klass med en behållare av andra
objekt.  Därtill förutsätts \emph{Behållare: Listor} och \emph{Behållare:
Tupler} för listor, skivning och tupler, \emph{Klasser och objekt} för
egna klasser med \mintinline{python}{__init__} och
\mintinline{python}{__str__}, \emph{Inmatning och felhantering} för
\mintinline{python}{try}/\mintinline{python}{except} och för att kasta
egna särfall, samt \emph{Moduler och paket} för att dela upp ett program
i moduler.

\emph{Läsning}
Uppgifterna hänvisar till Pythons egen dokumentation: kapitlet
\emph{Data Structures} i handledningen, avsnittet \emph{collections} i
biblioteksreferensen och avsnittet \emph{Emulating numeric types} i
språkreferensen.
